% Day 05：悬挂法找二维平面薄板的重心
\begin{tikzpicture}[scale=0.88,>=Latex,line cap=round,line join=round,
  every node/.style={font=\small}]

  % 同一薄板从孔 A 悬挂：把 A-C 连线转到竖直方向。
  \begin{scope}[shift={(-2.55,0.45)},rotate=-60.5]
    \coordinate (Cone) at (0,0);
    \coordinate (Aone) at (-0.92,0.52);
    \filldraw[fill=gray!12,draw=black!75,thick]
      (-1.10,-0.65) -- (0.90,-0.55) -- (1.20,0.40) --
      (0.20,1.00) -- (-1.00,0.55) -- cycle;
    \fill (Aone) circle (1.5pt);
  \end{scope}
  \draw[thick] ($(Aone)+(-0.22,0.16)$) -- ($(Aone)+(0.22,0.16)$);
  \draw[dashed,gray!70] (Aone) -- ($(Aone)+(0,-2.05)$);
  \draw[->,blue!75!black,thick] (Cone) -- ++(0,-0.72)
    node[right] {$G$};
  \node[above left] at (Aone) {$A$};
  \node[font=\footnotesize] at (-2.55,-1.20) {从孔 $A$ 悬挂};

  % 同一薄板从孔 B 悬挂：薄板转到另一稳定姿态。
  \begin{scope}[shift={(2.55,0.45)},rotate=121.4]
    \coordinate (Ctwo) at (0,0);
    \coordinate (Btwo) at (0.90,-0.55);
    \filldraw[fill=gray!12,draw=black!75,thick]
      (-1.10,-0.65) -- (0.90,-0.55) -- (1.20,0.40) --
      (0.20,1.00) -- (-1.00,0.55) -- cycle;
    \fill (Btwo) circle (1.5pt);
  \end{scope}
  \draw[thick] ($(Btwo)+(-0.22,0.16)$) -- ($(Btwo)+(0.22,0.16)$);
  \draw[dashed,gray!70] (Btwo) -- ($(Btwo)+(0,-2.05)$);
  \draw[->,blue!75!black,thick] (Ctwo) -- ++(0,-0.72)
    node[right] {$G$};
  \node[above right] at (Btwo) {$B$};
  \node[font=\footnotesize] at (2.55,-1.20) {从孔 $B$ 悬挂};

  \node[gray!75,font=\footnotesize] at (0,-1.58)
    {稳定时：重心在悬挂点正下方};

  % 把两次铅垂线转印到薄板自身，交点就是同一个重心。
  \begin{scope}[shift={(0,-3.35)}]
    \coordinate (C) at (0,0);
    \coordinate (A) at (-0.92,0.52);
    \coordinate (B) at (0.90,-0.55);
    \filldraw[fill=gray!12,draw=black!75,thick]
      (-1.10,-0.65) -- (0.90,-0.55) -- (1.20,0.40) --
      (0.20,1.00) -- (-1.00,0.55) -- cycle;
    \fill (A) circle (1.3pt);
    \fill (B) circle (1.3pt);
    \draw[dashed,red!75!black,thick] (A) -- (1.08,-0.61);
    \draw[dashed,red!75!black,thick] (B) -- (-1.00,0.61);
    \filldraw[fill=red!85!black,draw=red!85!black] (C) circle (2.2pt);
    \node[red!75!black,below=2pt] at (C) {重心};
  \end{scope}
  \node[align=center,font=\footnotesize] at (0,-4.45)
    {把两次铅垂线画回同一薄板，\\交点即为重心};
\end{tikzpicture}
